\subsection{Cause 6: Missing List and Type Unification Cases in the Substitution Engine}
\label{sec:andante-cause6}

While attempting to extend coverage to the \textsc{synthesis-sorted}
dataset (learning a predicate over Prolog list-structured examples), two
further, independent defects were identified in Andante's substitution
engine (\texttt{andante/substitution.py}), both pre-existing and
unrelated to the anti-unification issues described in
Section~\ref{sec:andante-cause5}.

\paragraph{Cause 6a: list terms were never wired into
\texttt{Substitution.unify()}.} Andante represents Prolog list terms
(e.g.\ \texttt{[H|T]}) with a dedicated \texttt{List} class, and this
class implements its own \texttt{unify()} method, used by the query
solver during proof search. However, \texttt{Substitution.unify()} --
the separate unification routine used specifically during bottom-clause
construction -- dispatches purely on \texttt{isinstance} checks against
\texttt{Constant}, \texttt{Variable}, \texttt{Function} and \texttt{Type},
and had no case at all for \texttt{List}. Consequently, any modeh or
modeb declaration with a list-typed argument (e.g.\
\texttt{modeh(1,f(+list))}) crashed with a \texttt{TypeError} as soon as
a ground example list was unified against its instantiated mode atom.
Similarly, \texttt{TreeShapedKnowledge.match()} and \texttt{.add()}
(\texttt{andante/knowledge.py}), used to index and retrieve background
clauses, indexed each clause argument by branching on
\texttt{Constant}/\texttt{Variable}/\texttt{Function} only; a list-typed
argument fell through every branch, leaving the corresponding position
unindexed and causing \texttt{set.intersection()} to later be called with
zero arguments -- an unrelated-looking \texttt{TypeError} whose true
cause was this missing case.

\emph{Fix.} \texttt{Substitution.unify()} was extended with explicit
\texttt{List}-\texttt{Variable} and \texttt{List}-\texttt{List} cases; the
latter delegates to \texttt{List.unify()} itself (reusing its already
correct handling of partial list patterns such as \texttt{[H|T]} matched
against a longer ground list), applying the resulting bindings directly
to the enclosing \texttt{Substitution} rather than reimplementing the
matching logic. \texttt{TreeShapedKnowledge.add()}, \texttt{.match()} and
\texttt{.remove()} were extended with a \texttt{List} case that treats a
list-typed argument as a wildcard for indexing purposes -- a safe
over-approximation, since any clause it admits as a candidate is still
subject to genuine unification afterward, so no false positives can
result, at the cost of losing fine-grained indexing specifically on
list-shaped arguments.

\paragraph{Cause 6b: unifying two \texttt{Type} objects had no
handler.} A second, independent defect was found in the same method:
\texttt{Substitution.unify()} handles the case of a \texttt{Type}
declaration (e.g.\ \texttt{+element}) unified against a
\texttt{Constant}, \texttt{Function} or \texttt{Variable}, but had no
branch for unifying two \texttt{Type} objects directly against each
other. This situation arises naturally whenever the same variable
appears twice as an argument of one literal under the same mode
declaration -- for instance \texttt{geq(V,V)} generalized against
\texttt{modeb(*,geq(+element,+element))}: the second occurrence of
\texttt{V} unifies the \texttt{Type("+","element")} already stored for
it against the freshly derived \texttt{Type("+","element")} for the same
argument slot. Because no \texttt{if}/\texttt{elif} branch matched this
case, the method fell through silently and returned \texttt{None}
instead of raising or succeeding. This \texttt{None} was then written
into the substitution's internal dictionary, and crashed a
\emph{subsequent, unrelated} call the next time any entry of that same
substitution was revisited, since \texttt{None} has no \texttt{.apply()}
method -- making the fault appear disconnected from its true cause and
substantially harder to trace.

\emph{Fix.} An explicit \texttt{Type}-\texttt{Type} branch was added:
if both \texttt{Type} objects agree in sign and name, unification
succeeds trivially (returning either one); otherwise a
\texttt{SubstitutionError} is raised, consistent with every other
branch of the method.

\paragraph{Validation.} With both fixes applied, \textsc{synthesis-sorted}
parses and runs to completion without error (previously it failed to
even construct a bottom clause). The resulting hypothesis search
correctly reaches literals chaining across the list structure (e.g.\
\texttt{head/2}, \texttt{tail/2} applied recursively), but converges to
an empty hypothesis: the target concept (a list being fully sorted)
requires comparing an unbounded number of consecutive element pairs,
which is inherently recursive and cannot be expressed by a single
non-recursive bottom clause of the kind Andante (and Progol) constructs
from one example. This is a representational limitation of the
formalism itself -- addressable only through predicate invention, which
Andante does not implement -- and not a further engine defect. No
regression was observed on any previously-validated dataset (family,
short\_family, trains, trains2, trains1000, zendo, zendo1,
iggp-rps), each of which continues to report identical hypotheses,
accuracy and running time to those measured prior to this fix.

\begin{table}[h]
\centering
\begin{tabular}{@{}p{3cm}p{5cm}p{5cm}@{}}
\toprule
Location & Symptom & Fix \\
\midrule
\texttt{substitution.py}, \texttt{knowledge.py} (Cause 6a) &
List-typed mode arguments crashed bottom-clause construction with a
\texttt{TypeError}, and background-knowledge indexing crashed with
\texttt{set.intersection()} called on zero arguments &
Added \texttt{List} cases to \texttt{Substitution.unify()} (delegating
to \texttt{List.unify()}) and to \texttt{TreeShapedKnowledge}'s
add/match/remove (treated as a safe wildcard for indexing) \\
\addlinespace
\texttt{substitution.py} (Cause 6b) &
Unifying two \texttt{Type} objects (e.g.\ from a repeated variable in
one literal) silently returned \texttt{None}, later crashing an
unrelated substitution lookup with an \texttt{AttributeError} &
Added an explicit \texttt{Type}-\texttt{Type} branch: succeeds if sign
and name match, otherwise raises \texttt{SubstitutionError} \\
\bottomrule
\end{tabular}
\caption{Summary of the Cause 6 fixes (list and type unification).}
\end{table}
